% ============================================================
%  Příprava na maturitu – Mechatronika
%  Kompiluj: pdflatex maturita_mechatronika.tex
%  Kódování: UTF-8
% ============================================================

\documentclass[a4paper,10pt]{article}

% --- Balíčky ---
\usepackage[T1]{fontenc}
\usepackage[utf8]{inputenc}
\usepackage[czech]{babel}
\usepackage[left=2.4cm, right=1.8cm, top=1.8cm, bottom=1.8cm]{geometry}
\usepackage{lmodern}
\usepackage{microtype}
\usepackage{parskip}
\usepackage{titlesec}
\usepackage{enumitem}
\usepackage{xcolor}
\usepackage{hyperref}
\usepackage{tabularx}
\usepackage{booktabs}
\usepackage{array}
\usepackage{tikz}
\usepackage{eso-pic}
\usepackage{mdframed}
\usepackage{amsmath}
\usepackage{amssymb}
\usepackage{pgfplots}
\pgfplotsset{compat=1.18}
\usetikzlibrary{arrows.meta, positioning, shapes.geometric, calc}

% --- Barvy ---
\definecolor{accent}{HTML}{03254E}
\definecolor{stripeblue}{HTML}{568EA3}
\definecolor{medgray}{HTML}{888888}
\definecolor{lightblue}{HTML}{EAF2F8}
\definecolor{lightyellow}{HTML}{FDFBE4}
\definecolor{lightgreen}{HTML}{EAF7EA}

% --- Modrý pruh vlevo ---
\AddToShipoutPictureBG{%
  \begin{tikzpicture}[remember picture, overlay]
    \fill[stripeblue]
      (current page.north west)
      rectangle
      ([xshift=10mm] current page.south west);
    \fill[accent!60]
      ([xshift=10mm] current page.north west)
      rectangle
      ([xshift=12mm] current page.south west);
  \end{tikzpicture}%
}

\hypersetup{colorlinks=true, urlcolor=accent, linkcolor=accent}

\titleformat{\section}
  {\large\bfseries\color{accent}}
  {\thesection.}
  {0.5em}
  {}
  [\color{accent}\titlerule]

\titlespacing{\section}{0pt}{16pt}{6pt}

\newmdenv[
  backgroundcolor=lightblue,
  linecolor=stripeblue,
  linewidth=1pt,
  roundcorner=4pt,
  innerleftmargin=8pt,
  innerrightmargin=8pt,
  innertopmargin=6pt,
  innerbottommargin=6pt,
  skipabove=4pt,
  skipbelow=4pt
]{pojmybox}

\newmdenv[
  backgroundcolor=lightyellow,
  linecolor=accent!40,
  linewidth=1pt,
  roundcorner=4pt,
  innerleftmargin=8pt,
  innerrightmargin=8pt,
  innertopmargin=6pt,
  innerbottommargin=6pt,
  skipabove=4pt,
  skipbelow=8pt
]{vysvetlenibox}

\newmdenv[
  backgroundcolor=lightgreen,
  linecolor=accent!30,
  linewidth=1pt,
  roundcorner=4pt,
  innerleftmargin=8pt,
  innerrightmargin=8pt,
  innertopmargin=6pt,
  innerbottommargin=6pt,
  skipabove=4pt,
  skipbelow=8pt
]{vzorcebox}

\pagestyle{plain}

% ============================================================
\begin{document}

% --- Záhlaví ---
\begin{minipage}[t]{0.75\textwidth}
    {\Huge\bfseries\color{accent} Příprava na maturitu}\\[4pt]
    {\large\color{stripeblue} Mechatronika}\\[2pt]
    \textcolor{medgray}{\small Suchánek Lukáš \quad SPŠ Prosek \quad 2026}
\end{minipage}
\hfill
\begin{minipage}[t]{0.22\textwidth}
    \raggedleft
    \begin{tikzpicture}
        \draw[color=stripeblue, rounded corners=4pt, line width=1pt]
            (0,0) rectangle (3,1.2)
            node[midway, align=center, color=accent, font=\small\bfseries]
            {Otázek: 22};
    \end{tikzpicture}
\end{minipage}

\vspace{6pt}
\noindent\color{accent}\rule{\linewidth}{1.5pt}
\vspace{2pt}

% ============================================================
\section{Senzorika -- teplota, síla, hmotnost, deformace, napětí, tlak}

\begin{pojmybox}
\textbf{Klíčové pojmy:}
termistor NTC/PTC, termoelektrický článek, RTD (Pt100),
IR teploměr, tenzometr, Wheatstoneův můstek, piezoelektrický senzor,
tlakový senzor, snímač síly
\end{pojmybox}

\begin{vysvetlenibox}
\textbf{Teplotní senzory}

\begin{center}
\begin{tabularx}{0.95\linewidth}{@{} l X l l @{}}
    \toprule
    \textbf{Typ} & \textbf{Princip} & \textbf{Rozsah} & \textbf{Přesnost} \\
    \midrule
    Pt100/Pt1000 & Odporová závislost platiny na teplotě & $-200$ až $+850$\,°C & $\pm0{,}1$\,°C \\
    Termočlánek & Seebeckův jev -- 2 různé kovy & $-270$ až $+2300$\,°C & $\pm1$\,°C \\
    NTC termistor & Záporný teplotní koeficient odporu & $-50$ až $+150$\,°C & $\pm0{,}5$\,°C \\
    IR bezdotykový & Záření černého tělesa & $-50$ až $+2000$\,°C & $\pm2$\,°C \\
    \bottomrule
\end{tabularx}
\end{center}

\begin{vzorcebox}
Pt100: $R(T) = R_0 \left(1 + AT + BT^2\right)$,
kde $R_0 = 100\,\Omega$ při $0\,°C$, $A = 3{,}9083\times10^{-3}$\,°C$^{-1}$\\
Seebeckův jev: $U = S_{AB} \cdot (T_1 - T_2)$,
kde $S_{AB}$ je Seebeckův koeficient [µV/°C]
\end{vzorcebox}

\textbf{Tenzometr a Wheatstoneův můstek}\\
Tenzometr mění odpor při mechanické deformaci:

\begin{vzorcebox}
$$GF = \frac{\Delta R / R}{\varepsilon}$$
kde $GF$ je gauge factor (typicky 2--2,5 pro kovové tenzometry),
$\varepsilon = \Delta l / l$ je relativní deformace.
\end{vzorcebox}

Tenzometry se zapojují do \textbf{Wheatstoneova můstku} pro eliminaci
teplotní chyby a zvýšení citlivosti:

\begin{vzorcebox}
$$U_{out} = U_{in} \cdot \frac{R_1 R_3 - R_2 R_4}{(R_1+R_2)(R_3+R_4)}$$
Plný můstek (4 aktivní tenzometry): 4× vyšší citlivost než čtvrtmůstek.
\end{vzorcebox}

\textbf{Tlakové senzory}
\begin{itemize}[noitemsep, leftmargin=1.4em, topsep=4pt]
    \item \textbf{Piezorezistivní} -- tenzometry na membráně;
          statický i dynamický tlak; nejčastější typ
    \item \textbf{Kapacitní} -- deformace membrány mění kapacitu;
          velmi přesné, vhodné pro nízké tlaky
    \item \textbf{Piezoelektrický} -- generuje náboj při tlaku;
          pouze dynamické měření (nárazy, hluk)
\end{itemize}
\end{vysvetlenibox}

% ============================================================

% ============================================================
\end{document}
